\clearpage
\section*{September 8, 2026: Production workflow review before merging}
\textit{Agent-drafted research entry; review of checkpoint 3f9e4a8.}

The production workflow does not yet meet the agreed research-workflow standard. The review inventoried 26 task Makefiles and 40 task-local R scripts outside audits, together with the shared helpers and document entry points. Successful builds against the current cache do not establish reliable rebuilding. No research estimates or sample definitions were changed by this review.

A disposable two-task example using the repository's shared Make rules reproduced a stale-result error under the installed GNU Make 3.81. After changing an upstream source, the first downstream build regenerated the upstream file but left the downstream result unchanged. A second invocation updated the downstream result. The existing-output rule uses an order-only prerequisite for its recursive check. Replacing that edge with a normal prerequisite in the disposable example updated the downstream result on the first invocation and left a subsequent unchanged build idle. That candidate correction has not been applied to the project or validated across its full graph.

The root empirical target has a separate omission: an existing figure-guide PDF has no prerequisite that checks its producing task. A dry run with the guide's LaTeX source marked changed scheduled nothing at the root, while the same check within the analysis task scheduled compilation. Multiple-output recovery also fails: the rule for a missing secondary output runs a file-size check rather than its producer. An isolated reproduction of that rule failed instead of restoring the file.

The source workflow has further gaps. Several acquisition and staging tasks declare only inventory CSVs, while scripts create and read the underlying releases outside concrete Make rules. Seven production Makefiles omit a sourced shared helper from their prerequisites. In a simulated transfer failure, the existing download helper deleted an already valid destination and returned a failure status instead of failing the build itself. No real source file was used or removed in this test.

Production still depends on the exposure classification and Attorney General evidence held under audits. The adopted classification and its required evidence need clear production ownership; manual review is not a reason to classify an input to the main results as exploratory. Conversely, the parent-distribution task still emits case listings and an explicitly named unit audit in production. Only the borough task currently declares standard task-local data reports. Shared helpers and execution rules are split between the task root, the old helper directory, and the new shared-code directory.

This review does not certify the branch as ready to merge. The priority is to repair concrete dependency and missing-output behavior, preserve failed-download safety, and then separate adopted data preparation from exploratory diagnostics. Shared-code placement, data reports, and source-file targets should be updated with their callers. Verification must include changed-input propagation and missing-output recovery, not only cached builds. The review did not rerun every historical source and audit or validate every numerical research assumption.

\textit{Reproduction:} The stale-result example copied the current
\texttt{tasks/generic.make} into a disposable two-task directory. The download test evaluated the existing download helper with a simulated failing transfer and a temporary destination. The root guide check used Make's dry-run and assumed-new-source options. All test files were outside the repository; the production code remains at the reviewed checkpoint.
